\documentclass[11pt]{article}
\usepackage[margin=1in]{geometry}
\usepackage{amsmath,amssymb,amsthm}
\usepackage[T1]{fontenc}
\usepackage{lmodern}
\usepackage{microtype}
\usepackage[hidelinks]{hyperref}
\newtheorem{proposition}{Proposition}
\newcommand{\qN}{q_N}
\title{A repeated interlace root below $-4$\\\large Disproof of conjecture 00000003474}
\author{gaochengzhecpu}
\date{October 4, 2026}
\begin{document}
\maketitle
\noindent\textbf{Result.} For the disjoint union $G=P_5\sqcup P_5$ of two
five-vertex paths, the original one-variable interlace polynomial is
\[
 \qN(G;x)=x^2(x^2+5x+2)^2.
\]
Consequently $r=(-5-\sqrt{17})/2<-4$ is a multiple real root. This contradicts
the conjecture's first clause, that all multiple roots lie in $[-4,0]$.

\paragraph{Definition and scope.}
We use the original single-variable interlace polynomial of
Arratia, Bollob\'as, and Sorkin~\cite{ABS}, subsequently called the
\emph{vertex-nullity interlace polynomial}. It is characterized on finite
simple graphs by
\begin{equation}\label{eq:rec}
 \qN(E_n;x)=x^n,\qquad
 \qN(H;x)=\qN(H-a;x)+\qN(H^{ab}-b;x)\quad(ab\in E(H)).
\end{equation}
Here $E_n$ is edgeless. For the ABS pivot $H^{ab}$, divide the other vertices
into those adjacent to $a$ only, $b$ only, both, or neither; toggle edges
between different classes among the first three. Incidences at $a,b$ remain
unchanged. The two source languages do not specify a different interlace
variant or restrict to connected graphs. The counterexample uses this
standard original convention, with no claim about a different polynomial.

\begin{proposition}
The graph $G=P_5\sqcup P_5$ has a multiple real interlace root outside $[-4,0]$.
\end{proposition}
\begin{proof}
Write $p_n=\qN(P_n;x)$, where $P_n$ has $n$ vertices and $p_0=1$.
Pivoting on an end edge $ab$, with $a$ the leaf, changes no edges.
Deleting $a$ leaves $P_{n-1}$; deleting $b$ leaves an isolated $a$ and
$P_{n-2}$. The polynomial is multiplicative on disjoint unions
\cite[Remark 16]{ABS}, giving $p_n=p_{n-1}+xp_{n-2}$ for $n\ge2$.
Starting with $p_1=x$ gives
\[
 p_2=2x,\quad p_3=x^2+2x,\quad
 p_4=3x^2+2x,\quad p_5=x^3+5x^2+2x.
\]
Multiplicativity now yields $\qN(G;x)=p_5^2=x^2(x^2+5x+2)^2$.
The quadratic has root $r=(-5-\sqrt{17})/2$. Since $17>9$, we have
$\sqrt{17}>3$ and hence $r<-4$. Its other root differs from $r$ and
$r\ne0$, so $r$ has multiplicity exactly two in $\qN(G;x)$.
One false necessary clause already disproves the full conjecture; its
additional embedding and link-spectrum clauses are not needed.
\end{proof}

\newpage
\paragraph{Lean verification.}
The accompanying Lean 4.19.0 project, with Mathlib pinned to commit
\texttt{c44e0c8ee63ca166450922a373c7409c5d26b00b}, contains:
\begin{itemize}
 \item actual Boolean adjacency matrices on $\mathrm{Fin}(5)$ and
 $\mathrm{Fin}(10)$, with symmetry and absence of loops checked;
 \item vertex deletion, the ABS pivot, and the total recursion
 \eqref{eq:rec}, choosing the first edge in lexicographic order;
 \item kernel-checked evaluation of that recursion on both graphs,
 proving the two polynomial identities above;
 \item proofs that the ten-vertex polynomial is nonzero, that
 $(X-r)^2$ divides it over $\mathbb R$, and that $r<-4$;
 \item a theorem negating the universal real multiple-root location claim
 for this interlace polynomial.
\end{itemize}
To keep finite graph evaluation small, an auxiliary recursion records the
orders of the edgeless leaves. A general Lean theorem proves that the sum
of their monomials equals the polynomial recursion. Ordinary
\texttt{decide} checks the finite reduction lists; polynomial algebra is
proved by kernel-checked tactics. There is no use of \texttt{native\_decide},
unfinished proof, or custom axiom.

\paragraph{Formalization boundary.}
The Lean definition directly implements the published original pivot
recursion. The general theorem that its value is independent of pivot
order, and its equivalence to the subset-nullity expression below, are
the established results cited in \cite{ABS,ABS2}; they are not re-proved
in this file. No hypothesis assigning a precomputed polynomial to $G$
is assumed. The formal multiple-root certificate proves multiplicity at
least two, which is all the disproof requires; the exact multiplicity
statement follows from the elementary factorization in the paper.

\paragraph{Independent exact computation.}
The pure-Python check enumerates all induced vertex subsets and computes
their adjacency-matrix ranks over $\mathbb F_2$, using
\[
 \qN(H;x)=\sum_{S\subseteq V(H)}(x-1)^{\dim\ker A_H[S]}.
\]
For $P_5$, the numbers of subsets of nullities $0,1,2,3$ are respectively
$8,15,8,1$. Expanding gives $x^3+5x^2+2x$. For $G$, all $1024$ subsets are
checked, giving ascending coefficients
$(0,0,4,20,29,10,1)$, exactly the square of the path polynomial.
This computation is supplementary and is not imported as a Lean axiom.

\begin{thebibliography}{9}\small
\bibitem{ABS} R.~Arratia, B.~Bollob\'as, and G.~B.~Sorkin,
\emph{The interlace polynomial of a graph},
J. Combin. Theory Ser. B \textbf{92} (2004), 199--233.
Definition 5, Theorem 12, and Remark 16.
\href{https://arxiv.org/abs/math/0209045}{arXiv:math/0209045}.
\bibitem{ABS2} R.~Arratia, B.~Bollob\'as, and G.~B.~Sorkin,
\emph{A two-variable interlace polynomial},
Combinatorica \textbf{24} (2004), 567--584, Section 4.
\href{https://arxiv.org/abs/math/0209054}{arXiv:math/0209054}.
\end{thebibliography}
\end{document}
